\chapter*{Abstract}
\addcontentsline{toc}{chapter}{Abstract}

This dissertation studies a shared computational question in two different settings: what can be learned by training a system to undo corruption? The first setting is sparse representation learning for natural images. The second is recurrent population coding of location. In both, corruption creates a supervised restoration signal from otherwise unlabeled or synthetically specified states. The objects being restored, however, are different--an image in one study and a recurrent activity pattern in the other--so the common framework is a method of analysis rather than a claim of identical biological implementation.

The first study proposes learning overlapping dependencies among sparse coefficients by differentiating an image-denoising loss through a finite sparse-inference procedure. The method recovers the intended plane-versus-axis structure of a three-dimensional toy distribution. On whitened natural-image patches, the learned groups combine localized, oriented filters with similar position, orientation, and spatial scale. The reported table gives higher patch output signal-to-noise ratio for the fully learned group model than for the compared factorial and fixed-topographic variants at each tested corruption level, although the margin over the strongest fixed-topographic model is small and no multi-seed uncertainty is available.

The second study trains recurrent dynamics, under state and weight perturbations, to write, hold, and transport prescribed one-dimensional location codes. A symmetric recurrent component is used for restoration, while a velocity-modulated antisymmetric component transports activity. In selected simulations with a fixed total neuron count, a mixed place--periodic code has lower location error than the place-only code in parts of the tested noise regime and exhibits a finer staircase of long-run fixed points. These findings establish a computational example, not a general explanation for biological grid cells: decoder details, model-selection protocol, budget matching beyond neuron count, multiple modules, and two-dimensional tests remain unresolved.

The synthesis is that restoration and transport should be separated. Denoising can identify directions that return a perturbed state toward a high-density or prescribed set, whereas transport moves a state along that set. This distinction makes the central thesis concrete, connects it to work on score estimation, learned iterative inference, continuous attractors, and neural population manifolds, and exposes the experiments needed for stronger claims.
